\daytitle{2}{质点、参考系、位移与速度}

\begin{learningbox}{本节导读}
从初中“路程”的概念进入高中“位移”的概念；讲清楚参考系为什么重要；区分平均速度和瞬时速度；建立描述运动的基本语言——位置、位移、速度，为后续分析速度变化打基础。
\end{learningbox}

\textbf{本节读题路线：} 先选参考系和坐标轴，再标出初位置、末位置和运动路径；接着判断题目问的是位移还是路程、速度还是速率。若出现相对运动或小船过河，先画速度矢量，再做矢量合成。

\section*{一、学习目标}

\begin{enumerate}
  \item 理解参考系的概念，知道同一个物体的运动在不同参考系下描述不同。
  \item 严格区分位移和路程：位移是矢量，路程是标量。
  \item 掌握平均速度和瞬时速度的区别与联系。
  \item 能从位移与时间的关系写出 $v = \dfrac{\Delta x}{\Delta t}$，并在匀速直线运动（$t_0=0$）中使用 $v=\dfrac{x-x_0}{t}$。
  \item 理解“速度”在高中物理中默认指瞬时速度，速率指速度大小。
  \item 初步会用速度矢量合成处理相对运动和小船过河模型。
\end{enumerate}

\section*{二、初高中衔接}

初中学过“速度 = 路程 $\div$ 时间”，这个公式在初中够用，但在高中有一个重要缺陷：\textbf{它分不清方向}。

看一个例子：
\begin{itemize}
  \item 小明从家向东走 $100\ \mathrm{m}$ 到学校，用了 $50\ \mathrm{s}$。初中的算法：$v = 100 / 50 = 2\ \mathrm{m/s}$。
  \item 放学后，小明从学校向西走 $100\ \mathrm{m}$ 回到家，也用了 $50\ \mathrm{s}$。初中的算法：还是 $2\ \mathrm{m/s}$。
\end{itemize}
两个过程速度数值相同，但方向相反。初中阶段常不区分方向，高中必须把方向写进物理量。这就是为什么高中要引入\textbf{位移}代替路程，用\textbf{矢量}代替标量。

另外，初中通常默认“大地是静止的”，高中阶段要进一步意识到：\textbf{静止是相对的}——这就是\textbf{参考系}。

\textbf{科学发展视角：} 导航定位、运动捕捉和自动驾驶都离不开参考系、位置、速度和相对运动。手机导航给出的“速度”来自位置随时间的变化，飞机和船舶导航还要同时考虑自身速度、风速或水流速度。本节的位移和速度，是理解现代定位与控制系统的第一步。

\section*{三、核心概念}

\begin{conceptbox}{参考系}

\textbf{直观理解：} 描述一个物体是运动还是静止，先得选一个“参照物”。坐在火车上观察窗外，会感觉树在向后运动——树相对于地面是静止的，但相对于火车是向后运动的。这个被选作标准的物体就是\textbf{参考系}。

\textbf{严格定义：} 在描述一个物体的运动时，被选作参考标准的另一个物体（或物体系），称为参考系。

\textbf{关键点：}
\begin{itemize}
  \item 同一个物体的运动，在不同参考系下描述不同。这叫\textbf{运动的相对性}。
  \item 高中物理中，如果不特别说明，默认以\textbf{地面}为参考系。
  \item 选择参考系的原则：使问题描述最简单。比如研究火车运动，选地面比选太阳方便得多。
\end{itemize}

\end{conceptbox}

\begin{conceptbox}{质点（回顾）}

Day 01 已经介绍过，质点是有质量但忽略大小形状的点。本节正式开始用质点模型来描述运动。

研究一个物体从 A 到 B 的运动时，只关心它的\textbf{位置变化}，不关心它自身如何旋转、是否变形。这时物体可以抽象为一个质点，用一条有向线段表示它的运动。

\begin{center}
% TikZ 图：位移与路程的对比
% Day 02：质点、参考系、位移与速度
% 使用：% TikZ 图：位移与路程的对比
% Day 02：质点、参考系、位移与速度
% 使用：% TikZ 图：位移与路程的对比
% Day 02：质点、参考系、位移与速度
% 使用：\input{figures/tikz/displacement-path.tex}
\begin{tikzpicture}[scale=1,>=Latex]
  % 坐标轴
  \draw[->] (-0.5,0) -- (6,0) node[right] {$x$};
  \draw[->] (0,-0.5) -- (0,2) node[above] {$y$};
  % 起点终点
  \filldraw[black] (1,0.8) circle (2pt) node[below] {$A(x_1,y_1)$};
  \filldraw[black] (4.5,1.2) circle (2pt) node[below] {$B(x_2,y_2)$};
  % 位移矢量
  \draw[->,thick,red] (1,0.8) -- (4.5,1.2) node[midway,above] {$\Delta \vect{r}$};
  % 轨迹虚线（路程）
  \draw[dashed] (1,0.8) .. controls (2,1.8) and (3.5,0.2) .. (4.5,1.2);
  \node[above] at (3,2.2) {实际路径（路程 $s$）};
  \node[red,below] at (2.8,0.4) {位移 $\Delta\vect{r}$};
\end{tikzpicture}

\begin{tikzpicture}[scale=1,>=Latex]
  % 坐标轴
  \draw[->] (-0.5,0) -- (6,0) node[right] {$x$};
  \draw[->] (0,-0.5) -- (0,2) node[above] {$y$};
  % 起点终点
  \filldraw[black] (1,0.8) circle (2pt) node[below] {$A(x_1,y_1)$};
  \filldraw[black] (4.5,1.2) circle (2pt) node[below] {$B(x_2,y_2)$};
  % 位移矢量
  \draw[->,thick,red] (1,0.8) -- (4.5,1.2) node[midway,above] {$\Delta \vect{r}$};
  % 轨迹虚线（路程）
  \draw[dashed] (1,0.8) .. controls (2,1.8) and (3.5,0.2) .. (4.5,1.2);
  \node[above] at (3,2.2) {实际路径（路程 $s$）};
  \node[red,below] at (2.8,0.4) {位移 $\Delta\vect{r}$};
\end{tikzpicture}

\begin{tikzpicture}[scale=1,>=Latex]
  % 坐标轴
  \draw[->] (-0.5,0) -- (6,0) node[right] {$x$};
  \draw[->] (0,-0.5) -- (0,2) node[above] {$y$};
  % 起点终点
  \filldraw[black] (1,0.8) circle (2pt) node[below] {$A(x_1,y_1)$};
  \filldraw[black] (4.5,1.2) circle (2pt) node[below] {$B(x_2,y_2)$};
  % 位移矢量
  \draw[->,thick,red] (1,0.8) -- (4.5,1.2) node[midway,above] {$\Delta \vect{r}$};
  % 轨迹虚线（路程）
  \draw[dashed] (1,0.8) .. controls (2,1.8) and (3.5,0.2) .. (4.5,1.2);
  \node[above] at (3,2.2) {实际路径（路程 $s$）};
  \node[red,below] at (2.8,0.4) {位移 $\Delta\vect{r}$};
\end{tikzpicture}


\figcaption{位移与路程的对比：位移是从 $A$ 指向 $B$ 的有向线段，路程是实际路径长度}\label{fig:displacement-path}
\end{center}

\end{conceptbox}

\begin{conceptbox}{位移}

\textbf{直观理解：} 位移就是“从起点指向终点的有向线段”。它只关心\textbf{起点和终点的位置差}，不关心中间走过什么路线。

\textbf{严格定义：} 位移 $\Delta \vect{x}$ 是描述物体位置变化（即位置矢量的变化）的物理量，是矢量。
\[
\Delta \vect{x} = \vect{x}_2 - \vect{x}_1
\]
其中 $\vect{x}_1$ 是初位置矢量，$\vect{x}_2$ 是末位置矢量。

\textbf{位移合成：} 位移是矢量，多个位移相加要遵循矢量加法。平行四边形法则的做法是：把两个位移矢量平移到同一起点，以它们为邻边作平行四边形，共同起点到对角点的对角线就是合位移。连续运动中也常用首尾相接的三角形法则，两种画法等价。

\textbf{位移 vs 路程：}

\begin{center}
\begin{tabular}{lll}
\toprule
 & 位移 $\Delta\vect{x}$ & 路程 $s$ \\
\midrule
性质 & 矢量 & 标量 \\
方向 & 有方向（初位置$\to$末位置） & 无方向 \\
大小 & 初末位置的直线距离 & 实际路径总长度 \\
关系 & $|\Delta\vect{x}| \leq$ 路程（直线不回头时相等） & $\geq$ 位移大小 \\
\bottomrule
\end{tabular}
\end{center}

\textbf{经典例子：} 绕操场跑一圈 $400\ \mathrm{m}$，位移 $= 0$（回到起点），路程 $= 400\ \mathrm{m}$。位移为零，路程不为零。

\end{conceptbox}

\begin{learningbox}{坐标语言中的运动学量分析清单}
高中运动学不是先记公式，而是先把物理量问清楚。沿一条直线建立坐标轴后，常见量可以这样读：

\begin{center}
\begin{tabular}{llll}
\toprule
物理量 & 回答的问题 & 数学表达 & 性质 \\
\midrule
位置 $x$ & 物体在哪里 & 坐标值 & 可正可负 \\
位移 $\Delta x$ & 从哪里到哪里 & $x_2-x_1$ & 矢量，有方向 \\
路程 $s$ & 实际走了多长 & 各段路径长度之和 & 标量，非负 \\
平均速度 $\bar v$ & 位置变化有多快 & $\dfrac{\Delta x}{\Delta t}$ & 矢量，有正负 \\
平均速率 & 路径走得有多快 & $\dfrac{s}{\Delta t}$ & 标量，非负 \\
\bottomrule
\end{tabular}
\end{center}

\textbf{解题顺序：} 先看题目给的是坐标还是路径长度；再用末位置减初位置得到位移；最后根据题目问“速度”还是“速率”选择分子。只要物体折返、绕圈或走曲线，就要特别警惕：路程通常不能代替位移。
\end{learningbox}

\begin{conceptbox}{速度}

\textbf{直观理解：} 速度描述物体运动快慢和方向。初中只讲“快慢”，高中还要讲“往哪走”。

\textbf{平均速度：}
\[
\vect{\bar{v}} = \frac{\Delta \vect{x}}{\Delta t}
\]
\begin{itemize}
  \item 矢量，方向与位移 $\Delta \vect{x}$ 相同。
  \item 它描述的是\textbf{一段时间内}整体的运动快慢和方向。
  \item 这里用的是位移 $\div$ 时间，不是路程 $\div$ 时间；后者叫\textbf{平均速率}（标量）。
\end{itemize}

\textbf{瞬时速度：}
\[
\vect{v} = \lim_{\Delta t \to 0} \frac{\Delta \vect{x}}{\Delta t}
\]
\begin{itemize}
  \item 矢量，描述物体在\textbf{某一时刻（或某一位置）}的运动快慢和方向。
  \item 高中物理中说“速度”，\textbf{默认指瞬时速度}。
  \item 瞬时速度的大小叫做\textbf{速率}（标量）。
\end{itemize}

\textbf{常见误用：}
\begin{enumerate}
  \item 平均速度 $\neq$ 速度的平均值。例如：某物体沿同一直线同向运动，前一半位移的速度为 $4\ \mathrm{m/s}$，后一半位移的速度为 $6\ \mathrm{m/s}$。速度的算术平均是 $(4+6)/2 = 5\ \mathrm{m/s}$，但全程平均速度要这样算：设每半程位移大小为 $s$，则总时间 $t = \dfrac{s}{4} + \dfrac{s}{6} = \dfrac{5s}{12}$，全程平均速度 $\bar{v} = \dfrac{2s}{5s/12} = \dfrac{24}{5} = 4.8\ \mathrm{m/s} \neq 5\ \mathrm{m/s}$。两段位移相等但时间不同，算术平均值法失效；应回到定义式 $\bar{v}=\Delta x/\Delta t$。
  \item 平均速率 $\neq$ 平均速度的大小。走弯路时平均速率大于平均速度的大小（因为路程 $\geq$ 位移大小）。
\end{enumerate}

\end{conceptbox}

\section*{四、公式与适用条件}

\textbf{公式使用前检查：}
\begin{enumerate}
  \item 先选参考系和坐标轴，说明正方向。
  \item 分清题目给的是位置、位移、路程，还是路径长度。
  \item 判断要求的是平均速度、瞬时速度，还是平均速率。
  \item 所有时间间隔都要用同一段过程，不能把不同过程的 $\Delta x$ 和 $\Delta t$ 混在一起。
\end{enumerate}

\textbf{核心公式：}

\begin{enumerate}
  \item \textbf{平均速度定义式：}
  \[
  \bar{v} = \frac{\Delta x}{\Delta t}
  \]
  \begin{itemize}
    \item 适用条件：任意一段运动都可用它定义该时间段的平均速度。
    \item 物理意义：一段时间内位置变化的平均快慢。
    \item 注意：这是\textbf{定义式}，不是决定式。它给出的是一段时间内的平均效果，不能直接替代某一时刻的瞬时速度。
  \end{itemize}

  \item \textbf{匀速直线运动的速度公式：}
  \[
  v = \frac{x-x_0}{t}\quad (t_0=0),\qquad
  \text{更一般地}\quad v=\frac{x_2-x_1}{t_2-t_1}
  \]
  \begin{itemize}
    \item 适用条件：物体做匀速直线运动（速度大小方向都不变）。
    \item 物理意义：单位时间内的位移。
    \item 注意：$x$ 是末位置坐标，$x_0$ 是初位置坐标，$x-x_0$ 才是位移。只有当计时起点 $t=0$ 时物体在原点 $x_0=0$，才可简写为 $v=x/t$。
    \item 常见误用：不能把匀加速运动的平均速度公式 $\bar{v} = (v_0+v)/2$ 与这个公式混淆。
  \end{itemize}

  \item \textbf{瞬时速度的极限定义：}
  \[
  v = \lim_{\Delta t \to 0} \frac{\Delta x}{\Delta t}
  \]
  \begin{itemize}
    \item 理解方式：把 $\Delta t$ 取得越来越小，趋近于零，平均速度就趋近于瞬时速度。
    \item 高中阶段不需要真的求极限，会用平均速度近似计算即可。
  \end{itemize}
\end{enumerate}

\section*{五、典型模型}

\textbf{直线运动模型：} 物体沿一条直线运动。这是高中物理最常见的模型，也是初学者最先掌握的。

在一维直线运动中，位移、速度的方向用\textbf{正负号}表示——设定一个正方向，与正方向同向为正，反向为负。

\begin{center}
% TikZ 图：一维直线运动坐标轴
% Day 02：质点、参考系、位移与速度
% 使用：% TikZ 图：一维直线运动坐标轴
% Day 02：质点、参考系、位移与速度
% 使用：% TikZ 图：一维直线运动坐标轴
% Day 02：质点、参考系、位移与速度
% 使用：\input{figures/tikz/1d-coordinate.tex}
\begin{tikzpicture}[scale=1,>=Latex]
  % 一维坐标轴
  \draw[->,thick] (-0.5,0) -- (7,0) node[right] {$x$};
  % 原点
  \node at (0,-0.3) {$O$};
  % 标记点
  \foreach \x/\label in {1/$x_1$,3/$x_2$,5/$x_3$} {
    \draw (\x,0.1) -- (\x,-0.1);
    \node at (\x,-0.3) {\label};
  }
  % 正方向箭头
  \node[above] at (6,0.15) {正方向 $\to$};
  % 位移箭头 x1->x2（正）
  \draw[->,red,thick] (1,0.3) -- (3,0.3) node[midway,above] {$\Delta x > 0$};
  % 位移箭头 x3->x2（负）
  \draw[->,red,thick] (5,0.65) -- (3,0.65) node[midway,above] {$\Delta x < 0$};
\end{tikzpicture}

\begin{tikzpicture}[scale=1,>=Latex]
  % 一维坐标轴
  \draw[->,thick] (-0.5,0) -- (7,0) node[right] {$x$};
  % 原点
  \node at (0,-0.3) {$O$};
  % 标记点
  \foreach \x/\label in {1/$x_1$,3/$x_2$,5/$x_3$} {
    \draw (\x,0.1) -- (\x,-0.1);
    \node at (\x,-0.3) {\label};
  }
  % 正方向箭头
  \node[above] at (6,0.15) {正方向 $\to$};
  % 位移箭头 x1->x2（正）
  \draw[->,red,thick] (1,0.3) -- (3,0.3) node[midway,above] {$\Delta x > 0$};
  % 位移箭头 x3->x2（负）
  \draw[->,red,thick] (5,0.65) -- (3,0.65) node[midway,above] {$\Delta x < 0$};
\end{tikzpicture}

\begin{tikzpicture}[scale=1,>=Latex]
  % 一维坐标轴
  \draw[->,thick] (-0.5,0) -- (7,0) node[right] {$x$};
  % 原点
  \node at (0,-0.3) {$O$};
  % 标记点
  \foreach \x/\label in {1/$x_1$,3/$x_2$,5/$x_3$} {
    \draw (\x,0.1) -- (\x,-0.1);
    \node at (\x,-0.3) {\label};
  }
  % 正方向箭头
  \node[above] at (6,0.15) {正方向 $\to$};
  % 位移箭头 x1->x2（正）
  \draw[->,red,thick] (1,0.3) -- (3,0.3) node[midway,above] {$\Delta x > 0$};
  % 位移箭头 x3->x2（负）
  \draw[->,red,thick] (5,0.65) -- (3,0.65) node[midway,above] {$\Delta x < 0$};
\end{tikzpicture}


\figcaption{一维直线运动的坐标轴与位移表示（正方向向右，$\Delta x > 0$ 向右，$\Delta x < 0$ 向左）}\label{fig:1d-coordinate}
\end{center}

在这个坐标系中，正方向向右。位移 $\Delta x > 0$ 表示位置变化向右，$\Delta x < 0$ 表示位置变化向左。

\nopagebreak[4]
\section*{六、例题精讲}

\begin{examplebox}{位移与路程的计算}

\textbf{题意分析：} 一个同学从校门口 $A$ 点出发，先向东走了 $300\ \mathrm{m}$ 到达 $B$ 点，然后向北走了 $400\ \mathrm{m}$ 到达 $C$ 点。
\begin{enumerate}
  \item 求该同学的总路程。
  \item 求该同学的位移大小和方向。
  \item 如果整个过程用时 $250\ \mathrm{s}$，求平均速度的大小和平均速率。
\end{enumerate}

\textbf{物理模型：} 质点在一维和二维平面上的运动模型。A$\to$B 是一维直线运动，再 A$\to$B$\to$C 是二维运动。

\textbf{解题过程：}

\begin{enumerate}
  \item 总路程 $s = AB + BC = 300 + 400 = 700\ \mathrm{m}$。
  
  \item 位移是从起点 $A$ 指向终点 $C$ 的有向线段。
  \[
  |\Delta\vect{x}| = \sqrt{AB^2 + BC^2} = \sqrt{300^2 + 400^2} = 500\ \mathrm{m}
  \]
  方向：东偏北 $\theta$ 角，其中 $\tan\theta = \dfrac{400}{300} = \dfrac{4}{3}$，$\theta \approx 53^\circ$。
  
  \item 平均速度大小：$|\bar{\vect{v}}| = \dfrac{|\Delta\vect{x}|}{\Delta t} = \dfrac{500}{250} = 2.0\ \mathrm{m/s}$，方向东偏北 $53^\circ$。
  
  平均速率：$\bar{v}_{\text{速率}} = \dfrac{s}{\Delta t} = \dfrac{700}{250} = 2.8\ \mathrm{m/s}$。
\end{enumerate}

\begin{center}
\begin{tikzpicture}[scale=0.78,>=Latex,line cap=round,line join=round,
  every node/.style={font=\small}]
  \coordinate (A) at (0,0);
  \coordinate (B) at (3,0);
  \coordinate (C) at (3,4);
  \coordinate (D) at (0,4);

  % Actual two-stage motion.
  \draw[->,very thick,blue!75!black] (A) -- (B)
    node[midway,below=3pt] {$300\ \mathrm{m}$};
  \draw[->,very thick,red!75!black] (B) -- (C)
    node[midway,right=3pt] {$400\ \mathrm{m}$};

  % Translated copies complete the parallelogram.
  \draw[->,dashed,red!70!black] (A) -- (D)
    node[midway,left=3pt] {$400\ \mathrm{m}$};
  \draw[->,dashed,blue!70!black] (D) -- (C);

  % Resultant displacement.
  \draw[->,ultra thick,orange!90!black] (A) -- (C)
    node[midway,above left=2pt] {$500\ \mathrm{m}$};

  \fill (A) circle (1.6pt) node[below left] {$A$};
  \fill (B) circle (1.6pt) node[below right] {$B$};
  \fill (C) circle (1.6pt) node[above right] {$C$};

  \draw (0.72,0) arc[start angle=0,end angle=53,radius=0.72];
  \node at (0.78,0.43) {$\theta$};

  % Direction key is separated from the displacement arrows.
  \begin{scope}[shift={(-1.05,-0.72)}]
    \draw[->] (0,0) -- (0.75,0) node[right] {东};
    \draw[->] (0,0) -- (0,0.75) node[above] {北};
  \end{scope}
\end{tikzpicture}


\figcaption{位移合成示例：平行四边形法则与三角形法则都得到合位移 $500\ \mathrm{m}$，方向东偏北 $\theta\approx 53^\circ$}\label{fig:displacement-composition}
\end{center}

\textbf{方法总结：}
\begin{enumerate}
  \item 路程就是实际路径长度之和，把每一段长度加起来即可。
  \item 位移用起点指向终点的直线段表示，二维情况下按平行四边形法则或三角形法则进行矢量合成。
  \item 平均速度 = 位移 $\div$ 时间，平均速率 = 路程 $\div$ 时间。两者不是同一概念。
  \item 画图辅助分析是解决运动学问题的基本功。
\end{enumerate}

\end{examplebox}

\begin{examplebox}{平均速度与速度的平均值}

\textbf{题意分析：} 某物体沿直线运动，前一半位移的平均速度为 $v_1 = 4\ \mathrm{m/s}$，后一半位移的平均速度为 $v_2 = 6\ \mathrm{m/s}$。求全程的平均速度。

\textbf{物理模型：} 直线运动，两段位移相等。

\textbf{解题过程：}

设总位移为 $2x$，前一半位移用时 $t_1$，后一半位移用时 $t_2$。
\[
t_1 = \frac{x}{v_1},\qquad t_2 = \frac{x}{v_2}
\]
全程平均速度：
\[
\bar{v} = \frac{2x}{t_1 + t_2} = \frac{2x}{\frac{x}{v_1} + \frac{x}{v_2}} = \frac{2}{\frac{1}{v_1} + \frac{1}{v_2}} = \frac{2v_1 v_2}{v_1 + v_2}
\]
代入 $v_1 = 4$，$v_2 = 6$：
\[
\bar{v} = \frac{2 \times 4 \times 6}{4 + 6} = \frac{48}{10} = 4.8\ \mathrm{m/s}
\]

\textbf{错解检查：} 有人会算 $(4 + 6)/2 = 5\ \mathrm{m/s}$，这个做法不成立，因为两段位移相等但速度不同，时间不相等，不能直接对速度取算术平均。

\textbf{方法总结：}
\begin{itemize}
  \item 平均速度最稳妥的入口是定义 $\bar{v} = \Delta x / \Delta t$，先求总位移和总时间。
  \item 等位移段的平均速度公式 $\bar{v} = \dfrac{2v_1v_2}{v_1+v_2}$ 是调和平均数。
  \item 类似地，等时间段的平均速度是算术平均数 $\bar{v} = \dfrac{v_1+v_2}{2}$。两种情形不要混淆。
\end{itemize}

\end{examplebox}

\section*{七、变式训练}

\begin{enumerate}
  \item 某人骑自行车沿直线从 $A$ 到 $B$，前 $1/3$ 路程的速度为 $4\ \mathrm{m/s}$，中间 $1/3$ 路程的速度为 $6\ \mathrm{m/s}$，最后 $1/3$ 路程的速度为 $12\ \mathrm{m/s}$。求全程的平均速度。（提示：不能直接对三段速度取算术平均。）
  
  \item 一质点沿半径为 $R = 10\ \mathrm{m}$ 的圆形轨道运动了半圈，用时 $5\ \mathrm{s}$。求：
  \begin{center}
  \begin{tikzpicture}[scale=0.82,>=Stealth]
  \coordinate (O) at (0,0);
  \coordinate (A) at (1.75,0);
  \coordinate (B) at (-1.75,0);
  \draw[gray!35,dashed] (O) circle (1.75);
  \draw[->,thick] (A) arc[start angle=0,end angle=180,radius=1.75];
  \draw[<->,blue,thick] (A) -- (B) node[midway,below=3pt] {位移大小 $=2R$};
  \draw[->] (O) -- (1.24,1.24) node[midway,above left] {$R$};
  \fill (A) circle (1.6pt) node[below right] {$A$};
  \fill (B) circle (1.6pt) node[below left] {$B$};
  \node at (0,2.13) {半圆路程 $=\pi R$};
\end{tikzpicture}

  \end{center}
  \begin{enumerate}
    \item 位移大小和路程。
    \item 平均速度大小和平均速率。
  \end{enumerate}
  
  \item 汽车从静止开始沿直线行驶，第 $1\ \mathrm{s}$ 内位移 $2\ \mathrm{m}$，第 $2\ \mathrm{s}$ 内位移 $4\ \mathrm{m}$，第 $3\ \mathrm{s}$ 内位移 $6\ \mathrm{m}$。求：
  \begin{enumerate}
    \item 前 $2\ \mathrm{s}$ 的平均速度。
    \item 第 $3\ \mathrm{s}$ 内的平均速度。
    \item 全程 $3\ \mathrm{s}$ 的平均速度。
  \end{enumerate}
\end{enumerate}

\section*{八、易错点}

\begin{mistakebox}
\begin{enumerate}
  \item \textbf{把路程当位移代入速度公式：} 高中速度公式 $v = \Delta x / \Delta t$ 中的分子是位移不是路程。题目问“速度”时必须先判断是平均速度还是平均速率。
  
  \item \textbf{平均速度和“速度的平均”混为一谈：} 平均速度是总位移 ÷ 总时间，不是各段速度的算术平均。等位移和等时间两种情况的公式不同，要会区分。
  
  \item \textbf{忽略参考系的选择：} 所有运动的描述都要基于某个参考系。题目不说明则以地面为参考系。但像“人在匀速行驶的船甲板上走”这类问题，人相对于船和相对于地的速度不同，必须明确参考系。
  
  \item \textbf{正负号只表示方向：} 在一维直线运动中，速度的正负只表示方向，不表示大小。$v = -10\ \mathrm{m/s}$ 和 $v = 5\ \mathrm{m/s}$ 相比，$|-10| > |5|$，前者更快。
  
  \item \textbf{混淆位移和路程：} 位移大小 $\leq$ 路程。单向直线运动时两者大小相等，但只要折返或走曲线，位移大小就小于路程。特别要注意“回到原点”时位移为零但路程不为零。

  \item \textbf{小船过河中把船头方向当作运动方向：} 有水流时，小船实际运动方向由船对水速度和水流速度的合速度决定。
\end{enumerate}
\end{mistakebox}

\section*{九、竞赛入门补充}

\begin{competitionbox}
\textbf{相对运动的概念}

在竞赛中，参考系的选取可以大大简化计算。一个经典技巧是：选择运动的物体作为参考系，把多个物体的运动转化为单个物体的运动。

\textbf{例题：} 一辆汽车以 $v_1 = 20\ \mathrm{m/s}$ 向东行驶，一辆摩托车以 $v_2 = 30\ \mathrm{m/s}$ 向东行驶。求摩托车相对于汽车的速度。

\textbf{解：} 取汽车为参考系。在汽车看来，地面以 $20\ \mathrm{m/s}$ 向西运动，摩托车相对于地面的速度是 $30\ \mathrm{m/s}$ 向东。因此摩托车相对于汽车的速度：
\[
\vect{v}_{2\text{对}1} = \vect{v}_2 - \vect{v}_1 = 30 - 20 = 10\ \mathrm{m/s}\ (\text{向东})
\]

更一般地，相对速度公式为：
\[
\vect{v}_{AB} = \vect{v}_A - \vect{v}_B
\]
其中 $\vect{v}_{AB}$ 表示 $A$ 相对于 $B$ 的速度。

\textbf{推广到二维：}
如果两辆车速度不共线，比如汽车向东 $20\ \mathrm{m/s}$，摩托车向北 $30\ \mathrm{m/s}$，则摩托车相对于汽车的速度：
\[
\vect{v}_{m\text{对}c} = \vect{v}_m - \vect{v}_c
\]
矢量减法用三角形法则计算，大小 $|\vect{v}_{m\text{对}c}| = \sqrt{20^2 + 30^2} \approx 36.1\ \mathrm{m/s}$，方向北偏西约 $33.7^\circ$。

这个思想在后续的追及相遇问题、抛体运动问题中会反复用到。

\medskip
\textbf{小船过河：速度合成的经典模型}

小船过河题本质上是相对速度的二维合成。若以地面为参考系：
\[
\vect{v}_{\text{船对地}}
=\vect{v}_{\text{船对水}}+\vect{v}_{\text{水对地}}.
\]
其中 $\vect{v}_{\text{船对水}}$ 由船头朝向和发动机决定，$\vect{v}_{\text{水对地}}$ 是水流速度，真正决定小船轨迹的是合速度 $\vect{v}_{\text{船对地}}$。

\begin{center}
% TikZ 图：小船过河的速度合成
% Day 02：质点、参考系、位移与速度
\begin{tikzpicture}[scale=0.86,>=Latex,line cap=round,line join=round,
  every node/.style={font=\small}]
  % Case 1: shortest time.
  \begin{scope}[shift={(-3.35,0)}]
    \fill[blue!5] (-1.85,-1.15) rectangle (1.85,1.15);
    \draw[thick] (-1.85,1.15) -- (1.85,1.15);
    \draw[thick] (-1.85,-1.15) -- (1.85,-1.15);
    \node[left] at (-1.85,1.15) {对岸};
    \node[left] at (-1.85,-1.15) {出发岸};
    \coordinate (O) at (-0.40,-0.75);
    \coordinate (W) at ($(O)+(0.90,0)$);
    \coordinate (S) at ($(O)+(0,1.50)$);
    \coordinate (R) at ($(O)+(0.90,1.50)$);
    \draw[->,blue!75!black,thick] (O) -- (W)
      node[midway,below] {$v_{\text{水}}$};
    \draw[->,red!80!black,thick] (O) -- (S)
      node[midway,left] {$v_{\text{船对水}}$};
    \draw[->,orange!85!black,thick] (O) -- (R)
      node[pos=0.62,right,xshift=1pt] {$v_{\text{船对地}}$};
    \draw[dashed,red!80!black] (W) -- (R);
    \draw[dashed,blue!75!black] (S) -- (R);
    \node[align=center] at (0,-1.70) {船头垂直河岸\\过河时间最短};
  \end{scope}

  % Case 2: shortest displacement.
  \begin{scope}[shift={(3.35,0)}]
    \fill[blue!5] (-1.85,-1.15) rectangle (1.85,1.15);
    \draw[thick] (-1.85,1.15) -- (1.85,1.15);
    \draw[thick] (-1.85,-1.15) -- (1.85,-1.15);
    \node[left] at (-1.85,1.15) {对岸};
    \node[left] at (-1.85,-1.15) {出发岸};
    \coordinate (P) at (-0.20,-0.75);
    \coordinate (W) at ($(P)+(0.72,0)$);
    \coordinate (S) at ($(P)+(-0.72,1.45)$);
    \coordinate (R) at ($(P)+(0,1.45)$);
    \draw[->,blue!75!black,thick] (P) -- (W)
      node[midway,below] {$v_{\text{水}}$};
    \draw[->,red!80!black,thick] (P) -- (S)
      node[midway,left] {$v_{\text{船对水}}$};
    \draw[->,orange!85!black,thick] (P) -- (R)
      node[midway,right] {$v_{\text{船对地}}$};
    \draw[dashed,red!80!black] (W) -- (R);
    \draw[dashed,blue!75!black] (S) -- (R);
    \draw[dashed,gray] (-0.20,-1.15) -- (-0.20,1.15);
    \node[align=center] at (0,-1.70) {抵消水流分量\\位移最短};
  \end{scope}
\end{tikzpicture}


\figcaption{小船过河：船对地速度由船对水速度与水流速度合成}\label{fig:boat-river-crossing}
\end{center}

\textbf{两个常见目标：}
\begin{enumerate}
  \item \textbf{过河时间最短：} 船头垂直河岸，让垂直河岸方向的速度最大。若河宽为 $d$，船相对水速度大小为 $v$，则 $t_{\min}=d/v$，但会被水流冲到下游。
  \item \textbf{位移最短：} 希望船相对地的合速度垂直河岸，就要让船头略向上游，使船速的上游分量抵消水流。若船速不足以抵消水流，则不能垂直到达对岸。
\end{enumerate}

\textbf{认知提醒：} “船头指向”不一定是“运动方向”。船头方向对应 $\vect{v}_{\text{船对水}}$，实际轨迹方向对应 $\vect{v}_{\text{船对地}}$。
\end{competitionbox}

\section*{十、今日小结}

\begin{enumerate}
  \item 参考系是描述运动的前提。运动是相对的，选择不同的参考系对同一运动的描述不同。
  \item 位移是矢量，从起点指向终点；路程是标量，是实际路径总长度。位移大小 $\leq$ 路程。
  \item 平均速度 $\bar{v} = \dfrac{\Delta x}{\Delta t}$，是矢量；平均速率 $= \dfrac{s}{\Delta t}$，是标量。
  \item 瞬时速度是 $\Delta t \to 0$ 时的极限，高中物理说“速度”默认指瞬时速度。
  \item 平均速度不是速度的平均值，必须用总位移 ÷ 总时间计算。
  \item 相对运动用速度矢量相加减；小船过河中，船对地速度等于船对水速度与水流速度的矢量和。
\end{enumerate}

\textbf{一分钟复盘：} 任选一个生活情境（电梯、公交车、小船过河），先说清参考系，再画位移或速度箭头，最后判断题目问的是位移、路程、平均速度还是平均速率。
